% 操作系统原理课内实践项目报告：题目五 连续内存分配与碎片问题验证
% 版式基线：Paper_Library/计算方法实验报告模板（汕头大学课程报告 LaTeX 模板库）
% 编译：XeLaTeX（中文必须 XeLaTeX）。VS Code 中 Ctrl+Alt+B 编译，Ctrl+Alt+V 预览。
\documentclass[12pt,a4paper]{ctexart}
% 中文默认使用 ctex 自带的 fandol 字体（TeX Live / MiKTeX 均自带）。

\usepackage{amsmath,amssymb,amsthm}
\usepackage{graphicx}
\usepackage{booktabs}
\usepackage{array}
\usepackage{tabularx}
\usepackage{geometry}
\usepackage{caption}
\usepackage{listings}
\usepackage{float}
\usepackage{enumitem}
\usepackage{xcolor}
\usepackage{tikz}
\usepackage{hyperref}
\hypersetup{hidelinks}

\geometry{left=2.6cm,right=2.6cm,top=2.6cm,bottom=2.6cm}
\graphicspath{{assets/},{../figures/}}
\captionsetup{font=small,labelsep=quad}
\renewcommand{\arraystretch}{1.15}

\lstset{
    language=Python,
    basicstyle=\ttfamily\footnotesize,
    keywordstyle=\bfseries,
    numbers=left,
    numberstyle=\tiny,
    numbersep=6pt,
    frame=single,
    breaklines=true,
    breakindent=2em,
    columns=fullflexible,
    keepspaces=true,
    showstringspaces=false,
    tabsize=4,
    xleftmargin=1.2em,
}

\renewcommand{\topfraction}{0.92}
\renewcommand{\bottomfraction}{0.85}
\renewcommand{\textfraction}{0.08}
\renewcommand{\floatpagefraction}{0.7}
\setcounter{topnumber}{3}
\setcounter{bottomnumber}{3}
\setcounter{totalnumber}{4}

\setlist{nosep,leftmargin=2em}

% ==================== 内存布局图命令 ====================
% \membar{总长度}{块列表}：块列表用逗号分隔，每项形如 类型/标签/长度
%   类型 f=空闲块, a=已分配块；标签取短名（地址由上方刻度标出）
%   例：\membar{90}{f/{空闲}/30, a/B/30, f/{空闲}/30}
\definecolor{freegreen}{RGB}{226,240,217}
\definecolor{allocblue}{RGB}{189,215,238}
\newcommand{\membar}[2]{%
\begin{tikzpicture}[x=\dimexpr\linewidth/#1\relax, y=0.9cm]
  \xdef\cursor{0}
  \foreach \btype/\blab/\blen [count=\bi] in {#2}{%
    \def\fftype{f}\def\aatype{a}%
    \ifx\btype\fftype
      \filldraw[fill=freegreen,draw=black] (\cursor,0) rectangle ({\cursor+\blen},1);
    \else
      \ifx\btype\aatype
        \filldraw[fill=allocblue,draw=black] (\cursor,0) rectangle ({\cursor+\blen},1);
      \fi
    \fi
    \pgfmathsetmacro{\midpt}{\cursor+\blen/2}%
    \node[font=\scriptsize] at (\midpt,0.5) {\blab};
    \ifnum\bi=1 \draw (\cursor,0) -- (\cursor,1.25)
      node[above,anchor=south west,font=\tiny]{\pgfmathprintnumber[fixed,precision=0]{\cursor}}; \fi
    \pgfmathparse{\cursor+\blen}\xdef\cursor{\pgfmathresult}
    \draw (\cursor,0) -- (\cursor,1.25)
      node[above,anchor=south east,font=\tiny]{\pgfmathprintnumber[fixed,precision=0]{\cursor}};
  }
\end{tikzpicture}\par}

% 代码附录在文件缺失时给出占位，保证模板可独立编译
\newcommand{\codefile}[1]{%
  \IfFileExists{#1}{\lstinputlisting{#1}}%
  {\itshape\small 待补充：实现完成后将 \texttt{#1} 放入仓库对应目录，本附录自动引入。\par}}

\begin{document}

% ==================== 封面 ====================
\begin{titlepage}
    \centering
    \vspace*{0.9cm}
    \includegraphics[width=3.8cm]{assets/logo.png}\\[0.9cm]
    {\zihao{2} \textbf{汕头大学}}\\[0.5cm]
    {\zihao{1} \textbf{操作系统原理课内实践项目报告}}\\[1.3cm]
    {\zihao{3} 题目五：连续内存分配与碎片问题验证}\\[1.8cm]

    \zihao{4}
    \renewcommand{\arraystretch}{1.55}
    \begin{tabular}{rl}
        组\hspace{2em}号： & 第 12 组 \\
        组\hspace{2em}长： & 陈子俊\ \ 2024612033 \\
        组\hspace{2em}员： & 黄应辉\ \ 2024611031 \\
                          & 蔡旭迪\ \ 2024612028 \\
                          & 史一飞\ \ 2024612031 \\
                          & 梁赫铭\ \ 2024612036 \\
                          & 李志恒\ \ 2024612066 \\
        指导教师： & 吴俊豪 \\
    \end{tabular}

    \vfill
    {\zihao{4} \the\year 年 \the\month 月 \the\day 日}
\end{titlepage}

% ==================== 目录（不编页码） ====================
\pagestyle{empty}
\tableofcontents
\clearpage

% ==================== 正文（页码从 1 开始） ====================
\pagenumbering{arabic}
\pagestyle{plain}

% ---------------------------------------------------------------
\section{模型构建}\label{sec:model}

\subsection{状态表示与模型约定}

设内存为大小固定 $T$ 的整数地址连续空间，任一时刻状态由已分配区与空闲区两组区间构成：
\begin{equation}
    \mathcal{A}=\{(id_i,\,[s_i,\,s_i+\ell_i))\},\qquad
    \mathcal{F}=\{[s_j,\,s_j+\ell_j)\}.
    \label{eq:state}
\end{equation}
初始 $\mathcal{A}=\varnothing$、$\mathcal{F}=\{[0,T)\}$。模型约定（不可检验，属规则设定）：
\textbf{H1} 整数地址、大小固定；\textbf{H2} 申请量为正整数，不计管理开销，不作内存紧缩；
\textbf{H3} 分配自选中空闲块的低地址端切出；\textbf{H4} 释放后立即与相邻空闲区合并。
每组序列不超过 20 步，申请使用唯一编号。每步记录四项指标：空闲总量、最大空闲块、空闲块数量、申请结果。

\subsection{两种分配策略}

申请量 $w$ 的候选集为 $\mathcal{C}(w)=\{F\in\mathcal{F}\mid \ell_F\ge w\}$，
为空则失败且状态不变。两种策略只在于选块规则不同：
\begin{equation}
    F_{\mathrm{FF}} = \arg\min_{F\in\mathcal{C}(w)} s_F ,
    \qquad
    F_{\mathrm{BF}} = \arg\min_{F\in\mathcal{C}(w)} \bigl(\ell_F,\ s_F\bigr).
    \label{eq:ff}
\end{equation}
即首次适应取地址最低的可用块，最佳适应取长度最小者（同长取低地址）。
选中 $[s,s+\ell)$ 后切出 $[s,s+w)$，余下 $[s+w,s+\ell)$ 仍为空闲块。

\subsection{释放与合并}

释放区 $[s,s+\ell)$ 的合并结果由两侧是否空闲决定：
\begin{equation}
\text{合并结果} =
\begin{cases}
[s_{L},\ s+\ell), & \text{左邻 } [s_{L},s)\in\mathcal{F} \\
[s,\ s_{R}), & \text{右邻 } [s+\ell,s_{R})\in\mathcal{F} \\
[s_{L},\ s_{R}), & \text{两侧均为空闲区} \\
[s,\ s+\ell), & \text{两侧均非空闲}
\end{cases}
\label{eq:merge}
\end{equation}
实现上将释放区按地址插入空闲表，随后自左向右扫描一次，凡首尾相接即合并。
该做法覆盖式 \eqref{eq:merge} 全部四种情形，且合并后仍继续向后检查，
可处理三块连续的链式合并。

\subsection{约束检查与实验假设}\label{subsec:checks}

每步自动检查三项：\textbf{(a)} 已分配区互不重叠；
\textbf{(b)} 守恒 $\sum_{\mathcal{A}}\ell+\sum_{\mathcal{F}}\ell=T$；
\textbf{(c)} 空闲表中无相邻块（合并完整）。
被拒绝的操作（失败申请、重复释放、释放不存在编号、申请量 $\le 0$ 或 $>T$）
必须报错且状态逐字段不变——否则守恒式立即被破坏，后续指标全部失真。

表 \ref{tab:hypo} 为六名成员各自独立提出的假设，提出时间早于程序运行与模型咨询。

\begin{table}[H]
    \centering
    \caption{六条待检验假设及其提出人与检验方式}
    \label{tab:hypo}
    \footnotesize
    \setlength{\tabcolsep}{4pt}
    \begin{tabularx}{\textwidth}{@{}c>{\raggedright\arraybackslash}X>{\raggedright\arraybackslash}X>{\raggedright\arraybackslash}p{1.7cm}>{\raggedright\arraybackslash}p{3.9cm}@{}}
        \toprule
        编号 & 假设内容 & 理由 & 提出人 & 检验方式 \\
        \midrule
        A1 & 申请成败只取决于最大空闲块，与空闲总量无关 & 分配要求地址连续 & 陈子俊 & 案例 3 第 6 步 \\
        A2 & 失败与选块策略无关，两策略在同一步失败 & 失败早于选块规则 & 黄应辉 & 案例 3 双策略轨迹比对 \\
        B2 & 最佳适应在所有负载下整体优于首次适应 & 每次剩余碎片最小 & 蔡旭迪 & 案例 1 分歧步 + 反例搜索 \\
        B5 & 策略差异经后续释放合并累积放大 & 合并依赖邻居状态 & 梁赫铭 & 案例 1 第 5--7 步 \\
        C1 & 关闭合并后空闲块数持续上升 & 每次释放新增一块 & 史一飞 & 合并开/关对照 \\
        D1 & 失败申请若切块会破坏守恒式 & 失败不是分配 & 李志恒 & 构造错误版本，用三项检查捕获 \\
        \bottomrule
    \end{tabularx}
\end{table}

% ---------------------------------------------------------------
\section{算法实现}\label{sec:algo}

\begin{sloppypar}
模拟器以状态对象（已分配表、空闲表）为核心，对外提供 \texttt{alloc(id,size)}、
\texttt{free(id)} 与 \texttt{check()}。主程序读入操作序列，逐步执行并输出四项指标与带地址的布局快照，
形成可复核的轨迹。两种策略共用同一执行框架与记录格式，仅替换式 \eqref{eq:ff} 的选择函数，
以保证对照公平。
\end{sloppypar}

首次适应按地址升序扫描、遇到第一个可用块即返回；最佳适应须遍历全部候选才能确定最小者，
两者复杂度同为 $O(|\mathcal{F}|)$，差异只在选择结果。合并实现的关键是
\textbf{合并后仍需继续向后检查}：若三块连续，第一次合并产生的新块还要与后一块再合并，
直到不再相邻；扫描结束时若仍存在相邻空闲区，即说明合并逻辑有遗漏，这同时构成了检查 (c) 的依据。

% ---------------------------------------------------------------
\section{实验设计}\label{sec:design}

表 \ref{tab:cases} 为三组序列。三组均从相同初始状态出发，分别以两种策略各运行一次；
序列在实验前定稿并冻结，输入与轨迹均纳入版本管理，保证报告中每条结论可追溯到具体步号。

\begin{table}[H]
    \centering
    \caption{三组实验序列的设计}
    \label{tab:cases}
    \small
    \setlength{\tabcolsep}{4pt}
    \begin{tabularx}{\textwidth}{@{}c>{\raggedright\arraybackslash}Xcc>{\raggedright\arraybackslash}X@{}}
        \toprule
        案例 & 设计目标 & 总内存 & 步数 & 预期现象 \\
        \midrule
        1 & 基本行为与策略分歧 & 100 & 7 & 第 5 步选块分歧，第 7 步合并类型不同 \\
        2 & 三种合并与非法操作 & 100 & 12 & 依次出现仅右、双侧、仅左合并；非法释放报错 \\
        3 & 外部碎片导致失败 & 90 & 9 & 第 6 步申请 50 失败：总量 60 够、最大块 30 不足 \\
        \bottomrule
    \end{tabularx}
\end{table}

案例 3 总内存取 90 而非 100，是为避免尾部残留空闲区干扰失败判定：三个 30 恰好占满内存，
释放首尾两块后必然形成"两小块夹一大块已分配区"的格局，使失败原因唯一、可手工验证。

人工核对取案例 3 做完整推演：逐笔手算每步布局与四项指标，与程序轨迹逐行对照。
人工推演只依赖式 \eqref{eq:ff}–\eqref{eq:merge} 的手工执行，不依赖代码正确性，
因此可作为独立证据；若两者不一致，则以第一个出现差异的步骤为界，判断是规则理解有误还是实现有缺陷。

% ---------------------------------------------------------------
\section{结果与分析}\label{sec:result}

\subsection{核心案例：外部碎片导致的申请失败}

表 \ref{tab:case3} 为案例 3 的逐步结果，图 \ref{fig:frag} 为失败瞬间的布局。
第 6 步申请 50 时，空闲总量已达 60，但两块各为 30，均小于 50。

\begin{table}[H]
    \centering
    \caption{案例 3 逐步轨迹（$T=90$，两策略结果一致）}
    \label{tab:case3}
    \small
    \setlength{\tabcolsep}{4pt}
    \begin{tabular}{cl>{\raggedright\arraybackslash}p{3.2cm}cccc}
        \toprule
        步 & 操作 & 布局（分配区$\mid$空闲区） & 空闲总量 & 最大块 & 结果 \\
        \midrule
        1--3 & alloc A/B/C 各 30 & A$[0,30)$ B$[30,60)$ C$[60,90)$$\mid$无空闲 & 0 & 0 & 成功 \\
        4 & free A & B、C 占用$\mid$空闲 $[0,30)$ & 30 & 30 & --- \\
        5 & free C & B 占用$\mid$空闲 $[0,30),[60,90)$ & 60 & 30 & --- \\
        6 & alloc D 50 & \textbf{不变} & 60 & 30 & \textbf{失败} \\
        7 & alloc E 20 & E$[0,20)$$\mid$空闲 $[20,30),[60,90)$ & 40 & 30 & 成功 \\
        8 & free B & 空闲 $[20,90)$ & 70 & 70 & --- \\
        9 & free E & 空闲 $[0,90)$ & 90 & 90 & --- \\
        \bottomrule
    \end{tabular}
\end{table}

\begin{figure}[H]
    \centering
    \membar{90}{f/{空闲}/30, a/{B}/30, f/{空闲}/30}
    \caption{案例 3 第 6 步失败瞬间的布局：空闲总量 60 足够，最大连续空闲块仅 30}
    \label{fig:frag}
\end{figure}

第 6 步支持 A2：失败发生在候选集为空之时，早于 $\arg\min$ 的选择，故两策略在同一步失败。
由此得到连续申请的判据：
\begin{equation}
    \text{申请 } w \text{ 成功} \iff \max_{F\in\mathcal{F}} \ell_F \ge w ,
    \label{eq:criterion}
\end{equation}
即空闲总量只是必要条件，\textbf{最大空闲块}才是判据（A1）。第 8、9 步显示系统并未受损——
释放 B 触发双侧合并、释放 E 触发右合并后完全恢复，
说明外部碎片是布局形态问题而非内存泄漏。

\subsection{两策略的分歧与累积效应}

表 \ref{tab:case1} 为案例 1 对照。第 1--4 步一致；第 5 步申请 10 时出现分歧：
首次适应取 $[0,30)$ 切出 $[0,10)$，最佳适应取恰好等大的 $[90,100)$ 整块用完（图 \ref{fig:case1ff}、\ref{fig:case1bf}）。

\begin{table}[H]
    \centering
    \caption{案例 1 双策略对照（$T=100$）}
    \label{tab:case1}
    \small
    \setlength{\tabcolsep}{4pt}
    \begin{tabular}{cl>{\raggedright\arraybackslash}p{4.2cm}>{\raggedright\arraybackslash}p{3.7cm}c}
        \toprule
        步 & 操作 & 首次适应 & 最佳适应 & 差异 \\
        \midrule
        1 & alloc A 30 & A$[0,30)$，空闲 $[30,100)$ & 同左 & \\
        2 & alloc B 30 & B$[30,60)$，空闲 $[60,100)$ & 同左 & \\
        3 & alloc C 30 & C$[60,90)$，空闲 $[90,100)$ & 同左 & \\
        4 & free A & 空闲 $[0,30),[90,100)$ & 同左 & \\
        5 & alloc D 10 & D$[0,10)$，空闲 $[10,30),[90,100)$ & D$[90,100)$，空闲 $[0,30)$ & $\star$ \\
        6 & free B & 左合并 $\to$ 空闲 $[10,60),[90,100)$ & 左合并 $\to$ 空闲 $[0,60)$ & \\
        7 & free C & 双侧合并 $\to$ 空闲 $[10,100)$ & 左合并 $\to$ 空闲 $[0,90)$ & \\
        \bottomrule
    \end{tabular}
\end{table}

\begin{figure}[H]
    \centering
    \membar{100}{a/{D}/10, f/{空闲}/20, a/{B}/30, f/{空闲}/30, f/{空闲}/10}
    \caption{案例 1 第 5 步后：首次适应把 D 放在低地址端}
    \label{fig:case1ff}
\end{figure}

\begin{figure}[H]
    \centering
    \membar{100}{f/{空闲}/30, a/{B}/30, f/{空闲}/30, a/{D}/10}
    \caption{案例 1 第 5 步后：最佳适应取走恰好等大的尾块}
    \label{fig:case1bf}
\end{figure}

第 6、7 步是同一串操作却在两策略下触发\textbf{不同合并类型}：首次适应第 7 步三块相邻、双侧合并，
最佳适应右侧邻居已分给 D，只能左合并。一次选块分歧由此改变整体布局的演化路径（B5）。
这也说明策略评价必须针对具体负载，离开负载谈"哪种更好"没有意义（B2 的反例基础）。

\subsection{三种合并型式}

\begin{table}[H]
    \centering
    \caption{三种合并型式的触发位置}
    \label{tab:merge}
    \small
    \setlength{\tabcolsep}{5pt}
    \begin{tabular}{ll>{\raggedright\arraybackslash}p{4.9cm}l}
        \toprule
        型式 & 触发步骤 & 合并前 & 合并后 \\
        \midrule
        仅右侧 & 案例 2 第 9、12 步 & 被释放区与右侧空闲区 & 合成一块 \\
        双侧 & 案例 2 第 7 步、案例 1 第 7 步 & 左侧、被释放区、右侧 & 三块合一 \\
        仅左侧 & 案例 2 第 6 步 & 左侧空闲区与被释放区 & 合成一块 \\
        \bottomrule
    \end{tabular}
\end{table}

\begin{figure}[H]
    \centering
    \membar{100}{f/{空闲}/20, a/{B}/20, f/{空闲}/20, a/{D}/20, f/{空闲}/20}

    \smallskip
    $\Downarrow$ \ 释放 B：左右两侧均为空闲区（案例 2 第 7 步）

    \smallskip
    \membar{100}{f/{空闲}/60, a/{D}/20, f/{空闲}/20}
    \caption{双侧合并（三块合一）前后对比}
    \label{fig:merge}
\end{figure}

三块合一是最易写错的情形：若只处理左合并或右合并之一，会留下两个相邻空闲区，违反检查 (c)。

\subsection{约束与非法操作检查}

\begin{table}[H]
    \centering
    \caption{约束检查与非法操作测试结果}
    \label{tab:checks}
    \small
    \begin{tabularx}{\textwidth}{@{}>{\raggedright\arraybackslash}X>{\raggedright\arraybackslash}Xc@{}}
        \toprule
        检查项 & 方法 & 结果 \\
        \midrule
        已分配区不重叠 & 逐步自动检查 + 轨迹抽样 & 通过 \\
        守恒 $\sum$已分配 $+\ \sum$空闲 $=T$ & 逐步自动检查 & 通过 \\
        合并完整（无相邻空闲区） & 逐步自动检查 & 通过 \\
        失败申请不改变状态 & 操作前后逐字段比对 & 通过 \\
        重复释放、释放不存在编号 & 构造非法输入 & 报错且不变 \\
        申请量 $\le 0$ 或 $>T$ & 构造非法输入 & 报错且不变 \\
        \bottomrule
    \end{tabularx}
\end{table}

故意构造的"失败也切块"版本被守恒检查立即捕获，说明三项检查对实现缺陷有效（D1）。
关闭合并的对照版本中，空闲块数随释放次数单调上升（C1）。

% ---------------------------------------------------------------
\section{结论}\label{sec:concl}

\begin{enumerate}
    \item 连续分配要求地址连续：申请 $w$ 成功的充要条件是最大空闲块不小于 $w$（式 \eqref{eq:criterion}），
          空闲总量只是必要条件（A1 成立）。
    \item 失败与选块策略无关：失败发生在候选集为空之时，两策略在同一序列上的失败步骤完全相同（A2 成立）。
    \item 最佳适应的"最佳"限于单次选择，它优先消耗恰好够用的块，中小申请反复出现时会把大块逐步碎化，
          故"整体优于首次适应"不能成立为一般结论（B2 被推翻）。
    \item 策略差异具有累积效应：一次选块分歧经后续合并放大为整体布局差别（B5 成立）。
    \item 合并可降低空闲块数量，但不能消除外部碎片（C1 成立）；失败与非法操作必须保持状态不变（D1 成立）。
\end{enumerate}

局限：模型不计管理开销、不作紧缩，结论适用于"整数地址、无重定位"的连续分配简化模型，
不能直接外推到真实系统（后者常用分页、伙伴系统等机制）；实验为三组不超过 20 步的序列，
属有限案例，只能否定某些判断，不能证明某策略普遍最优。

% ---------------------------------------------------------------
\appendix
\newpage
\section{程序代码}

\begin{sloppypar}
全部代码以 JavaScript 实现，运行于浏览器（无第三方依赖），可打开 \texttt{web/index.html} 本地复现。
\end{sloppypar}

\subsection*{A.1\quad 状态与检查（core.js）}
\codefile{../web/core.js}

\subsection*{A.2\quad 两种策略与合并（strategy.js）}
\codefile{../web/strategy.js}

\subsection*{A.3\quad 演示页面与轨迹输出（app.js）}
\codefile{../web/app.js}

\newpage
\section{小组分工与贡献}\label{app:division}

\begin{table}[H]
    \centering
    \caption{小组分工、对应假设与个人证据位置}
    \label{tab:division}
    \small
    \setlength{\tabcolsep}{4pt}
    \begin{tabularx}{\textwidth}{@{}llc>{\raggedright\arraybackslash}X>{\raggedright\arraybackslash}X@{}}
        \toprule
        成员 & 学号 & 假设 & 主要职责 & 个人成果与证据位置 \\
        \midrule
        陈子俊（组长） & 2024612033 & A1 & 内存模型与约束定义、报告整合 & 第 \ref{sec:model} 节约定、守恒与重叠检查 \\
        黄应辉 & 2024611031 & A2 & 首次适应实现 & 式 \eqref{eq:ff} 实现、双策略失败比对 \\
        蔡旭迪 & 2024612028 & B2 & 最佳适应实现 & 式 \eqref{eq:ff} 实现、同长规则核对、反例搜索 \\
        史一飞 & 2024612031 & C1 & 释放与合并测试 & 三种合并案例、合并开/关对照 \\
        梁赫铭 & 2024612036 & B5 & 演示页面与比较分析 & 布局与指标页面、三组输入、分歧步对照 \\
        李志恒 & 2024612066 & D1 & 约束检查与反例实验 & 三项自动检查、非法输入测试、故障版本对照 \\
        \bottomrule
    \end{tabularx}
\end{table}

项目按六人组织。相比任务书的五人分工，另设"约束检查与反例实验"一项，
承担三项逐步检查、非法输入测试与故障版本对照，并负责组内交叉复核。

\section{AI 工具使用说明}\label{app:ai}

\begin{sloppypar}
大语言模型仅用于辅助讨论与实现审查，完整问答与核验记录见仓库 \texttt{model-interactions/}。
\end{sloppypar}

\begin{table}[H]
    \centering
    \caption{大模型使用范围与学生判断}
    \label{tab:ai}
    \small
    \begin{tabularx}{\textwidth}{@{}>{\raggedright\arraybackslash}p{3.4cm}>{\raggedright\arraybackslash}X>{\raggedright\arraybackslash}X@{}}
        \toprule
        用途 & 涉及内容 & 核验方式与结论 \\
        \midrule
        策略比较与序列构造 & 请模型比较两策略并提出体现差异的序列 & 运行验证模型预测，据轨迹决定采纳或修改 \\
        失败原因解释 & 请模型解释"总量足够仍失败"及合并规则 & 用程序状态核对其是否混淆空闲总量与最大空闲块 \\
        绝对化判断审查 & 请模型评价"最佳适应整体更优" & 构造反例或说明适用条件后决定取舍 \\
        实现思路审查 & 请模型审查三种合并的处理顺序 & 用三块合一的边界案例验证 \\
        \bottomrule
    \end{tabularx}
\end{table}

以下内容由成员独立完成，未使用 AI 生成：第 \ref{sec:model} 节的状态模型与式 \eqref{eq:ff}–\eqref{eq:merge}
的整理与手算；核心算法实现（选块、切割、合并、三项检查）；表 \ref{tab:case1}--\ref{tab:checks} 的全部实验结果
均由本组程序运行产生，可按附录 A 复现。模型给出的序列与判断均经实际运行或手工推演核验，未直接采信其结论。

\end{document}
